\section{Introduction}

For $u>1$, the Dirichlet beta function is
\[
 \beta(u)=\sum_{k=0}^{\infty}\frac{(-1)^k}{(2k+1)^u}
         =1-\frac1{3^u}+\frac1{5^u}-\cdots.
\]
Its values at positive odd integers are nonzero rational multiples
of powers of $\pi$, and hence transcendental. The arithmetic of
its even values is much less understood.

For an even positive integer $s$, write
\[
 I(s)=\#\{j:1\le j\le s/2,\ \beta(2j)\notin\Q\}
\]
for the number of irrational values among
$\beta(2),\beta(4),\ldots,\beta(s)$.

In \cite{Tind} we proved that the rational span of
$\beta(2),\beta(4),\ldots,\beta(s)$ has dimension at least
$0.21\sqrt{s/\log s}$ for sufficiently large even $s$.
The same lower bound holds for $I(s)$; see Appendix~\ref{app:earlier-count}.

Here we show that an irrationality bound of the same shape and
a better constant is valid from a small explicit value of $s$.
In particular, Theorem~\ref{thm:main} gives a stronger count for
every even $s\ge160$. We use the constant\footnote{Replacing $c$ by
$0.358$ gives a bound of the same shape for every even $s\ge64$,
giving one irrational value at $64$, but only two rather than three
at $364$. For $64\le s\le158$, this follows from $I(64)\ge1$
(Section~\ref{sec:endpoints}) and
$0.358\sqrt{158/\log158}<2$; for $s\ge160$, use
Theorem~\ref{thm:main}.}
\[
 c=\pi\sqrt{\frac{2}{135}}=0.3823824806363651\ldots.
\]
All logarithms in this paper are natural.

\begin{theorem}\label{thm:main}
For every even integer $s\ge160$,
\begin{equation}\label{eq:main}
 I(s)\ge\left\lfloor c\sqrt{\frac{s}{\log s}}\right\rfloor.
\end{equation}
\end{theorem}
For every even $s\ge42\,170$, we will also prove the bound without the
floor.

The construction gives more precise information at individual
endpoints.  For each column below, at least $B$ of the even beta
values through $\beta(s)$ are irrational. For comparison, the last row
shows the expression inside the floor in \eqref{eq:main}, rounded to
two decimal places.
\begin{center}
\setlength{\tabcolsep}{5pt}
\begin{tabular}{@{}rrrrrrrrrr@{}}
\toprule
$B$&2&3&4&5&6&7&8&9&10\\
\midrule
$s$&160&364&588&1030&1296&1992&2502&3260&3828\\
\addlinespace[2pt]
$c\sqrt{s/\log s}$&2.15&3.00&3.67&4.66&5.14&6.19&6.84&7.68&8.24\\
\bottomrule
\end{tabular}
\end{center}
At $s=160$ and $s=364$, the expression inside the floor is
$2.147\ldots$ and $3.004\ldots$ respectively. Hence the uniform
bound~\eqref{eq:main} matches the counts in the table at these endpoints.
At $s=588$, the uniform bound gives $I(588)\ge3$ after taking the
floor of $3.67\ldots$, whereas applying the construction with
parameters chosen for this endpoint gives $I(588)\ge4$.

\subsection{A finite endpoint for every requested count}

We prove the main theorem by first giving an endpoint for every
positive integer $B$. Let $\varphi$ be Euler's totient function,
with $\varphi(1)=1$.

\begin{theorem}\label{thm:family}
For every integer $B\ge1$, put
\begin{equation}\label{eq:family}
 D_B=\sum_{d=1}^B\varphi(d),\qquad
 s_B=2\left\lceil D_B\left(33+\frac{45}{2}\log B\right)\right\rceil.
\end{equation}
Then $I(s_B)\ge B$. The sharper endpoints in the preceding table
also hold.
\end{theorem}

The quantity $D_B$ counts the distinct reduced fraction classes needed
for the $B$ grids in our construction. It grows quadratically in $B$;
the endpoint has size $B^2\log B$. Inverting this relation produces
the square root in Theorem~\ref{thm:main}.
The formula for $c$ is also of interest in its own right.
The $\pi$ and the factor $3$ in $135=3\cdot45$ come from
$D_B\sim(3/\pi^2)B^2$, the asymptotic count of reduced fractions.
The factor $45$ comes from balancing
denominator growth against decay, and the factor $2$ from
inverting the quadratic growth in $B$.
Section~\ref{sec:count}
makes the inversion explicit, including the constant and the small
endpoints.

\subsection{Related constructions}

Lai and Yu~\cite{LY20} obtained lower bounds of order
$\sqrt{s/\log s}$ for the number of irrational odd zeta values,
using rational functions whose zeros are arranged on rational grids.
Kumar and Tosi~\cite[Theorem~1]{KT26} adapt this approach to real
Dirichlet characters. For the beta function, their stated asymptotic
coefficient is $1.1510775\ldots$, with the starting endpoint depending
on the chosen loss in that coefficient. Theorem~\ref{thm:main} gives
a fixed bound at every even endpoint from $160$ onward.

The rational-function construction and coefficient elimination used
here build on these methods. The separation of coefficients by an
auxiliary prime has an antecedent in Lai and Sprang~\cite{LS25}.
A positive Gram matrix supplies the needed rank for every $B$. We give the
denominator, decay and prime-dilation arguments in full.

\subsection{How the contradiction will arise}

Fix $B$ and an even endpoint $s$. Suppose that fewer than $B$ of
the listed beta values are irrational. We construct $B$ sequences
of linear forms, indexed by positive odd integers $n$, of the shape
\begin{equation}\label{eq:preview}
 F_{n,b}=q_b(n)+\sum_{\substack{2\le i\le s\\i\text{ even}}}
                  b^i\rho_i(n)\beta(i),\qquad 1\le b\le B,
\end{equation}
with rational coefficients. The coefficient column at $\beta(i)$
is a scalar multiple of the fixed vector
$(1^i,2^i,\ldots,B^i)^{\mathsf T}$. Fixed integer weights on the
$B$ forms can therefore cancel any $B-1$ specified beta terms
at every index $n$. Choosing these terms to include the possible
irrational values produces a sequence of rational numbers.

For example, when $B=3$, the combination
$15F_{n,1}-6F_{n,2}+F_{n,3}$ cancels the coefficients at
$\beta(2)$ and $\beta(4)$, since
$15-6\cdot2^i+3^i=0$ for $i=2,4$. The same weights work for
every $n$ because the factor $\rho_i(n)$ is shared by the forms.

We bound the growth of a common denominator for the coefficients
and choose the parameters so that the forms decay faster than this
denominator grows. The rational beta values that remain after
cancellation contribute only a fixed additional denominator.
Thus the fixed combination, after clearing all denominators, gives
a sequence of integers tending to zero. To obtain a contradiction,
we must show that infinitely many of these integers are nonzero.

The cancellation weights can be chosen with a nonzero last entry.
Indeed, the first $B-1$ coordinates of the chosen power vectors
form a nonsingular generalized Vandermonde matrix
(Section~\ref{sec:contradiction}). We therefore arrange the constant
vector at the first index to have only its last coordinate nonzero:
\begin{equation}\label{eq:target}
 (q_1(1),\ldots,q_B(1))=\eta e_B,\qquad \eta\ne0,
\end{equation}
where $e_B=(0,\ldots,0,1)$. The combined form then has a nonzero
constant term at $n=1$. A congruence at indices $n=p$, for arbitrarily large primes
$p$, transfers this nonvanishing: after a suitable scaling, the
constant term survives modulo $p$ and all remaining beta terms
vanish. We obtain nonzero integers tending to zero, which is
impossible.

The main construction is the choice of \eqref{eq:target}. We find
$B$ polynomial multipliers whose constant vectors form a rational
invertible matrix. Its entries are integrals of products of
independent polynomials against one positive weight. This proves
invertibility for every $B$ and allows us to prescribe the seed
constant vector. The remaining sections check that these multipliers
preserve the denominator and decay estimates needed for the
contradiction.
